% Part: first-order-logic
% Chapter: axiomatic-deduction
% Section: proving-things-quant

\documentclass[../../../include/open-logic-section]{subfiles}

\begin{document}

\olfileid[fr]{fol}{axd}{prq}

\olsection{\usetoken{P}{derivation} avec quantificateurs}

\begin{ex}
Donnons une !!{derivation} de
$(\lforall[x][!A(x)] \land \lforall[y][!B(y)]) \lif
\lforall[x][(!A(x) \land !B(x))]$.

Remarquons d'abord que
\begin{align*}
  (\lforall[x][!A(x)] \land \lforall[y][!B(y)]) & \lif \lforall[x][!A(x)]\\
  \intertext{est une instance de l'\olref[prp]{ax:land1}, tandis que}
  \lforall[x][!A(x)] & \lif !A(a) \\
  \intertext{est une instance de l'\olref[qua]{ax:q1}. Ainsi, d'après la \olref[pro]{prop:chain},}
  (\lforall[x][!A(x)] \land \lforall[y][!B(y)]) & \lif !A(a) \\
  \intertext{est !!{derivable}. De même, puisque}
  (\lforall[x][!A(x)] \land \lforall[y][!B(y)]) & \lif \lforall[y][!B(y)] \qquad\text{et}\\
    \lforall[y][!B(y)] & \lif !B(a)\\
    \intertext{sont respectivement des instances de l'\olref[prp]{ax:land2} et de l'\olref[qua]{ax:q1},}
    (\lforall[x][!A(x)] \land \lforall[y][!B(y)]) & \lif !B(a)\\
    \intertext{est !!{derivable} d'après la \olref[pro]{prop:chain}. Une instance appropriée de l'\olref[prp]{ax:land3} et deux applications de~\MP{} montrent alors que}
    (\lforall[x][!A(x)] \land \lforall[y][!B(y)]) & \lif (!A(a) \land !B(a))\\
    \intertext{est !!{derivable}. Nous pouvons enfin appliquer \QR{} et obtenir}
(\lforall[x][!A(x)] \land \lforall[y][!B(y)]) & \lif \lforall[x][(!A(x) \land !B(x))].
\end{align*}
\end{ex}

\end{document}
